\documentclass[11pt]{amsart}
\usepackage[T1]{fontenc}
\usepackage[utf8]{inputenc}
\usepackage[french]{babel}
\usepackage{amssymb}
\usepackage[margin=1.05in]{geometry}
\usepackage[hidelinks]{hyperref}
\newtheorem{proposition}{Proposition}
\newtheorem{lemma}[proposition]{Lemme}
\theoremstyle{remark}
\newtheorem{remark}[proposition]{Remarque}
\DeclareMathOperator{\im}{im}
\newcommand{\Rees}{\mathcal R}
\newcommand{\F}{\mathcal F}
\title{Un certificat fini pour relever une section}
\date{26 octobre 2026}
\begin{document}
\maketitle

\noindent\textbf{Statut.}
Cette note donne une preuve algébrique constructive du relèvement
\emph{d'une section qui se relève à chaque ordre fini}, sous une hypothèse
explicite de finitude d'un module de cohomologie gradué.
La note complémentaire \texttt{finite-cohomology.tex} construit cette
finitude dans le cas considéré, par élimination et calcul de Čech.
Ces constructions restent à formaliser dans Unit. La présente note
ne prétend pas fournir une borne uniforme valable pour toutes les sections.

\section{Le certificat à vérifier}

Soient $A$ un anneau commutatif, $I$ un idéal, et
\[
 C^0\xrightarrow{d^0}C^1\xrightarrow{d^1}C^2,
 \qquad d^1d^0=0,
\]
un complexe de $A$-modules. On écrit $d=d^0$ et $C_n^j=I^nC^j$.
Les $C^j$ ne sont pas supposés de type fini sur $A$.

\begin{lemma}[Correction finie]\label{lem:correction}
Fixons $N\geq1$, $z_1,\ldots,z_N\in C^0$ avec
$z_i-z_1\in IC^0$. Supposons donnés
\[
 r_i\in I^{N-i}\quad(1\leq i<N),\qquad v\in I^NC^0,
 \qquad
 dz_N-\sum_{i=1}^{N-1}r_i\,dz_i=dv.
 \tag{*}
\]
Alors
\[
 z=z_N-\sum_{i=1}^{N-1}r_i z_i-v
 \quad\text{satisfait}\quad dz=0,\qquad z-z_1\in IC^0.
\]
\end{lemma}
\begin{proof}
La première égalité est $(*)$ et la linéarité de $d$.
Pour la seconde, $N-i\geq1$, donc tous les $r_i$ sont dans $I$,
et $v\in I^NC^0\subseteq IC^0$.
\end{proof}

Pour $I=(a,b)$, l'appartenance à $I^m$ se certifie sous la forme
\[
 r=\sum_{j=0}^{m}a^{m-j}b^j c_j.
\]
L'appartenance de $v$ à $I^NC^0$ se certifie de la même manière,
avec des coefficients dans $C^0$. Le certificat est donc fini :
un entier $N$, des cochaînes, des coefficients et une identité linéaire.
Sur un recouvrement affine explicite, on peut ensuite représenter les
cochaînes par des fractions et vérifier les identités avec des témoins
d'annulation après multiplication par des puissances des dénominateurs.

\section{Pourquoi une telle relation finit par apparaître}

Posons
\[
 \Rees=\bigoplus_{n\geq0} I^nT^n\subseteq A[T],
 \qquad
 M_n=H^1(C_n^\bullet)
   =\frac{\ker(d^1:C_n^1\to C_n^2)}{d^0(C_n^0)},
 \qquad M=\bigoplus_{n\geq0}M_n.
\]
La multiplication par $r\in I^j$ définit $M_n\to M_{n+j}$,
ce qui fait de $M$ un $\Rees$-module gradué.

\begin{proposition}\label{prop:lifting}
Supposons que $M$ possède la propriété noethérienne constructive
suivante : sur toute suite d'éléments de $M$, un terme est combinaison
linéaire des précédents, avec un témoin fini de cette relation.
Supposons donnés, pour tout $n\geq1$, des $z_n\in C^0$ tels que
\[
 dz_n\in I^nC^1,\qquad z_n-z_1\in IC^0.
 \tag{**}
\]
Alors il existe $z\in C^0$ avec $dz=0$ et $z-z_1\in IC^0$.
Seul un segment fini des données $(**)$ est utilisé.
\end{proposition}
\begin{proof}
La classe $\beta_n=[dz_n]$ appartient à $M_n$, puisque $d^1dz_n=0$.
Appliquons la propriété noethérienne à $\beta_1,\beta_2,\ldots$.
Elle fournit $N\geq1$ et une relation
\[
 \beta_N=\sum_{i=1}^{N-1}q_i\beta_i,\qquad q_i\in\Rees.
\]
Prendre la composante homogène de degré $N$ de cette relation donne
\[
 \beta_N=\sum_{i=1}^{N-1}(r_iT^{N-i})\beta_i,
 \qquad r_i\in I^{N-i}.
\]
La définition du quotient $M_N$ fournit $v\in C_N^0$ tel que
\[
 dz_N-\sum_{i=1}^{N-1}r_i\,dz_i=dv.
\]
On applique le lemme~\ref{lem:correction}.
\end{proof}

\begin{remark}
La compatibilité modulo toutes les puissances n'est pas nécessaire :
il suffit que les $z_n$ représentent la même classe modulo $I$.
Lorsque la propriété noethérienne est fournie par un arbre, comme dans
Unit, des relèvements donnés seulement sous troncation suffisent : on
parcourt cet arbre en choisissant à chaque étape un relèvement et son
code, puis on conclut sous troncation. C'est une récursion sur un arbre
bien fondé, et non un appel au choix dénombrable.
\end{remark}

\section{Le lien précis avec la noethérianité de Unit}

Dans l'application, $A=\mathbf Q[a,b]$ et $I=(a,b)$, donc
\[
 \Rees\simeq A[U,V]/(bU-aV),\qquad U\mapsto aT,\quad V\mapsto bT.
\]
Supposons fournis des générateurs $m_1,\ldots,m_r$ de $M$ sur
$\Rees$, avec la propriété de génération. Considérons l'anneau
à idéal de carré nul
\[
 E=\Rees\oplus M,\qquad
 (r,m)(s,n)=(rs,rn+sm).
\]
Il est engendré comme $\mathbf Q$-algèbre par les générateurs de
$\Rees$ et les $(0,m_j)$. La noethérianité constructive des algèbres
de type fini sur un corps discret s'y applique.
Une relation d'idéal
\[
 (0,\beta_N)=\sum_{i<N}(q_i,w_i)(0,\beta_i)
\]
est exactement une relation $\beta_N=\sum_{i<N}q_i\beta_i$.
Ainsi, une liste de générateurs de $M$ suffit : il n'est pas nécessaire
de calculer une présentation du module de tous les noyaux.

Les fichiers existants pertinents sont
\texttt{sag/noetherian/Generated.seg} (noethérianité par adjonctions
finies) et \texttt{sag/noetherian/Sequence.seg} (utilisation d'une
relation finie, parcours de l'arbre avec des représentants seulement
tronqués). L'anneau $\Rees\oplus M$ et les adaptations aux modules
et aux degrés restent à écrire. La proposition~\ref{prop:lifting}
est une variante graduée du mécanisme employé par \texttt{Sequence.Kill},
pas un résultat déjà exporté par ce fichier.

\section{Application au recouvrement projectif}

Soient $p,q\in A[x,y]$ tels que
\[
 p\equiv x\pmod{I},\qquad q\equiv y\pmod{I}.
\]
On note $W\subseteq\mathbf P^2_A$ le schéma fermé défini par leurs
homogénéisés et $\F=\mathcal O_W(-1)$.
Prenons le complexe de Čech des trois cartes standard pour $\F$.
Comme les cartes et leurs intersections sont affines, les sections
de $I^n\F$ sont les $I^nC^j$, et les sections de $\F/I^n\F$
se représentent par des $z_n\in C^0$ avec $dz_n\in I^nC^1$.
Ces assertions utilisent l'exactitude des sections quasi-cohérentes
sur les affines, pas la platitude sur $A$.

Ici, on peut choisir \emph{la même cochaîne à tous les ordres} :
\[
 z_*=(0,0,Z^{-1})\in C^0,\qquad
 dz_*=(0,Z^{-1},Z^{-1})\in C^1,
\]
les composantes doubles étant ordonnées $XY,XZ,YZ$.
En effet, écrivons $p=x+af+bg$ et $q=y+au+bv$.
Dans $B=A[x,y]/(p,q)$, on a $x=-af-bg$, $y=-au-bv$.
Sur l'intersection $XZ$, identifiée à $B[1/x]$, on obtient
\[
 1=a(-f/x)+b(-g/x).
\]
Donc $IB[1/x]=B[1/x]$, et de même $IB[1/y]=B[1/y]$ sur $YZ$.
En élevant ces identités à la puissance $n$, on construit les témoins
de $dz_*\in I^nC^1$ pour chaque $n$; la composante $XY$ est nulle.
Les hypothèses $(**)$ sont donc vérifiées avec $z_n=z_*$.
Cela évite ici toute construction séparée de relèvements de Hensel
ou d'idempotents. Les sections $h_n$ sont simplement les classes de
$z_*$ modulo $I^nC^0$.

La proposition~\ref{prop:lifting} s'applique dès que l'on a construit
la finitude sur $\Rees$ de
\[
 \bigoplus_{n\geq0}H^1(W,I^n\F).
 \tag{F}
\]
Elle donne $h\in H^0(W,\F)$ dont la restriction modulo $I$
est celle de $h_1$. Alors $Zh,Xh,Yh$ sont des sections globales
de $\mathcal O_W$. Si l'on a également construit la finitude de
$H^0(W,\mathcal O_W)$ sur $A$, leurs restrictions $s,sx,sy$
à la carte affine sont entières sur $A$, et $s\equiv1\pmod{IB}$,
où $B=A[x,y]/(p,q)$. Plus précisément, la correction s'écrit
$h=(1-\sum_{i<N}r_i)z_*-v$ et sa composante affine donne
$s=1-\sum_{i<N}r_i-Zv_Z$.

\section{Portée et travail restant}

Le module gradué de faisceaux $\bigoplus I^n\F$ est de type fini
sur l'algèbre de Rees tirée sur $W$. La finitude cohomologique
projective sur la base $\Rees$ donne (F); c'est le contenu classique
du lemme~30.20.1 des Stacks~\cite{stacks}. La note complémentaire
\texttt{finite-cohomology.tex} donne une construction des générateurs
dans le cas présent : une élimination fournit l'idéal de présentation
du module de Rees, puis chaque générateur homogène de degré projectif
$d$ fournit $\binom d2$ cocycles. Elle construit aussi les générateurs
des sections globales nécessaires au passage aux éléments entiers.
Ces arguments, ainsi que le relèvement décrit ici, restent à formaliser
dans Unit.

En revanche, une fois (F) donnée, le relèvement de la section voulue
ne demande ni les générateurs de
$\bigoplus_n\ker(H^1(I^n\F)\to H^1(\F))$, ni une borne uniforme
de stabilisation. Il utilise une relation sur la seule suite des
obstructions effectivement produites.
Il prouve l'implication suivante : une section modulo $I$ qui possède
des relèvements à chaque ordre fini possède un relèvement global
ayant la même réduction modulo $I$.
Il ne prouve pas l'énoncé plus fort donnant un seul ordre $N$, fixé
à l'avance, qui fonctionne pour toutes les sections.

Enfin, les termes du complexe de Čech ne sont généralement pas
finis sur $A$. On ne peut donc pas leur appliquer directement
Artin--Rees en les traitant comme des modules finis. De même,
un modèle fini quasi-isomorphe ne peut pas être réduit modulo $I^n$
sans contrôler sa compatibilité avec cette réduction.

\begin{thebibliography}{9}
\bibitem{stacks} The Stacks Project, section 30.20, en particulier
les lemmes 30.20.1 et 30.20.3--4,
\url{https://stacks.math.columbia.edu/tag/02O7}.
\end{thebibliography}
\end{document}
